\textbf{Proof.}
By ass:exposure-positivity, \(\Pi\) is invertible.  If \(B\in\mathcal B_Z\), the defining nondegeneracy condition gives \(S_B=B^{\top}\Omega B\succ0\), so \(S_B^{-1}\) exists.  Put
\[
E=\begin{bmatrix}0&I_n\\ I_n&0\end{bmatrix},
\qquad
A_B=\Pi^{-1}\Omega B S_B^{-1}B^{\top}\Omega\Pi^{-1}.
\]
The moment-matrix construction gives \(Q(c)=Q(0)-cE\).  Cyclicity and linearity of trace yield
\[
g_B(c)=n^{-2}\operatorname{tr}\{A_BQ(0)\}
-c n^{-2}\operatorname{tr}(A_BE),
\]
so \(g_B\) is affine.  The identity proved in lem:ht-affine-drift gives \(h(c)=h(0)+2c/n\).  Since def:gain-loss gives \(\ell_B(c)=h(c)-g_B(c)\), \(\ell_B\) is affine as well.  Ass:known-exposure and ass:iid-potential-pairs supply the causal exposure and ex ante moment interpretations; the affine calculation itself uses only the stated matrix identities, with all inverses justified above.

By ass:frechet-covariance-domain and def:coupling-class,
\[
\left\{\operatorname{Cov}_H\{Y_i(1),Y_i(0)\}:H\in\Gamma(F_1,F_0)\right\}
=[c_L,c_U].
\]
An affine real function on this compact interval attains its minimum and maximum at an endpoint.  Equality of the covariance image with the entire closed interval also says that each endpoint is attained by a coupling in \(\Gamma(F_1,F_0)\).  Hence
\[
\inf_{H\in\Gamma(F_1,F_0)}
g_B\!\left(\operatorname{Cov}_H\{Y_i(1),Y_i(0)\}\right)
=\min\{g_B(c_L),g_B(c_U)\}
\]
and
\[
\sup_{H\in\Gamma(F_1,F_0)}
\ell_B\!\left(\operatorname{Cov}_H\{Y_i(1),Y_i(0)\}\right)
=\max\{\ell_B(c_L),\ell_B(c_U)\},
\]
with both extrema attained.

By ass:product-coupling-zero, \(F_1\otimes F_0\in\Gamma(F_1,F_0)\) has covariance zero.  Thus ass:frechet-covariance-domain gives \(0\in[c_L,c_U]\), or \(c_L\le0\le c_U\), and \(g_B(0)\) is ordinarily typed.

Finally write \(g_B(c)=\alpha_B+\beta_Bc\).  If \(\beta_B=0\), then \(g_B\) is constant and its value at zero is its worst-coupling value.  Conversely suppose \(c_L<0<c_U\).  If \(\beta_B>0\), then \(g_B(c_L)=g_B(0)+\beta_Bc_L<g_B(0)\); if \(\beta_B<0\), then \(g_B(c_U)=g_B(0)+\beta_Bc_U<g_B(0)\).  Thus equality with the infimum is possible only when \(\beta_B=0\). \(\square\)